\documentclass{article}
\usepackage{amsmath,amssymb,amsthm}
\usepackage{algorithm}
\usepackage{algpseudocode}
\usepackage{geometry}
\geometry{margin=1in}
\newtheorem{theorem}{Theorem}
\newtheorem{lemma}{Lemma}
\newcommand{\grad}{\nabla}
\newcommand{\E}{\mathbb{E}}

\begin{document}

%%%%%%%%%%%%%%%%%%%%%%%%%%%%%%%%%%%%%%%%%%%%%%%%%%%%%%%%%%%%
\section*{Algorithm Name}
\textbf{Nesterov’s Accelerated Gradient with Polynomial Momentum}\\
(β‑schedule $\displaystyle \beta_k=\frac{k-1}{k+2}$)

%%%%%%%%%%%%%%%%%%%%%%%%%%%%%%%%%%%%%%%%%%%%%%%%%%%%%%%%%%%%
\section*{Mathematical Formulation}
Let $f:\mathbb{R}^d\rightarrow\mathbb{R}$ be the objective.
Initialize $x_{0}=y_{0}\in\mathbb{R}^d$ and choose a stepsize $\alpha>0$.
For $k=0,1,2,\dots$ perform
\[
\boxed{
\begin{aligned}
\beta_k &:= \frac{k-1}{k+2},\qquad (\beta_0:=0)\\[2mm]
y_k    &= x_k + \beta_k\,(x_k-x_{k-1}) ,\\[2mm]
x_{k+1}&= y_k - \alpha\,\grad f(y_k).
\end{aligned}}
\tag{1}
\]

%%%%%%%%%%%%%%%%%%%%%%%%%%%%%%%%%%%%%%%%%%%%%%%%%%%%%%%%%%%%
\section*{Pseudocode}
\begin{algorithm}[H]
\caption{Nesterov Accelerated Gradient (NAG) with $\beta_k=\frac{k-1}{k+2}$}
\begin{algorithmic}[1]
\State \textbf{Input:} stepsize $\alpha>0$, initial point $x_0$, maximum iterations $K$
\State $x_{-1}\gets x_0$, $x_0\gets x_0$
\For{$k=0$ \textbf{to} $K-1$}
    \State $\beta_k \gets \dfrac{k-1}{k+2}$ \Comment{Momentum coefficient}
    \State $y_k \gets x_k + \beta_k\,(x_k - x_{k-1})$
    \State $g_k \gets \grad f(y_k)$
    \State $x_{k+1} \gets y_k - \alpha\, g_k$
\EndFor
\State \textbf{Return} $x_{K}$
\end{algorithmic}
\end{algorithm}

%%%%%%%%%%%%%%%%%%%%%%%%%%%%%%%%%%%%%%%%%%%%%%%%%%%%%%%%%%%%
\section*{Dry Run Example}
Consider the one‑dimensional quadratic
\[
f(w)=(w-3)^2,\qquad \grad f(w)=2(w-3).
\]
Set $w_0=0$, stepsize $\alpha=0.1$, and run three iterations.

\begin{center}
\begin{tabular}{c|c|c|c|c}
$k$ & $\beta_k$ & $y_k$ & $g_k=\grad f(y_k)$ & $w_{k+1}=y_k-\alpha g_k$ \\ \hline
0 & $0$ & $0+0(0-0)=0$ & $2(0-3)=-6$ & $0-0.1(-6)=0.6$ \\[2mm]
1 & $\dfrac{0}{3}=0$ & $0.6+0(0.6-0)=0.6$ & $2(0.6-3)=-4.8$ & $0.6-0.1(-4.8)=1.08$ \\[2mm]
2 & $\dfrac{1}{4}=0.25$ & $1.08+0.25(1.08-0.6)=1.245$ & $2(1.245-3)=-3.51$ & $1.245-0.1(-3.51)=1.596$ 
\end{tabular}
\end{center}
Objective values:
\[
f(0)=9,\quad f(0.6)=5.76,\quad f(1.08)=3.6864,\quad f(1.596)=1.9729.
\]
The function value decreases rapidly, illustrating the acceleration effect.

%%%%%%%%%%%%%%%%%%%%%%%%%%%%%%%%%%%%%%%%%%%%%%%%%%%%%%%%%%%%
\section*{Assumptions}
\begin{enumerate}
\item \textbf{Convexity.} $f$ is convex:
\[
f(\theta x+(1-\theta)y)\le \theta f(x)+(1-\theta)f(y),\quad\forall x,y,\;\theta\in[0,1].
\tag{A1}
\]
\item \textbf{L–smoothness.} There exists $L>0$ such that
\[
\|\grad f(x)-\grad f(y)\|\le L\|x-y\|,\qquad\forall x,y\in\mathbb{R}^d.
\tag{A2}
\]
Equivalently,
\[
f(y)\le f(x)+\langle\grad f(x),y-x\rangle+\frac{L}{2}\|y-x\|^{2}.
\tag{A2'}
\]
\item \textbf{Stepsize.} The stepsize is chosen as 
\[
\alpha =\frac{1}{L}.
\tag{A3}
\]
\item \textbf{Bounded initial distance.} Define $R:=\|x_{0}-x^{\star}\|$, where $x^{\star}$ is a minimizer of $f$.
\end{enumerate}

%%%%%%%%%%%%%%%%%%%%%%%%%%%%%%%%%%%%%%%%%%%%%%%%%%%%%%%%%%%%
\section*{Main Convergence Theorem}
\begin{theorem}[Optimal $O(1/k^{2})$ rate]
\label{thm:main}
Under assumptions (A1)–(A3) the iterates generated by (1) satisfy for all $k\ge 0$
\[
f(x_{k})-f(x^{\star})\;\le\;\frac{2L\,R^{2}}{(k+2)^{2}}.
\tag{T1}
\]
Consequently $f(x_{k})\to f(x^{\star})$ with rate $\mathcal{O}(1/k^{2})$.
\end{theorem}

%%%%%%%%%%%%%%%%%%%%%%%%%%%%%%%%%%%%%%%%%%%%%%%%%%%%%%%%%%%%
\section*{Theoretical Properties}
\begin{itemize}
\item \textbf{Stability.} The momentum coefficient $\beta_k\in[0,1)$ for all $k$, guaranteeing that the extrapolation $y_k$ never overshoots the previous two iterates.
\item \textbf{Hyperparameter Sensitivity.} The schedule $\beta_k=(k-1)/(k+2)$ is derived from the optimal choice for convex smooth problems; any deviation destroys the $1/k^{2}$ guarantee.
\item \textbf{Comparison.} Standard (non‑accelerated) gradient descent with stepsize $1/L$ yields $f(x_k)-f(x^{\star})\le \frac{L R^{2}}{2k}$, i.e. a slower $\mathcal{O}(1/k)$ rate. The present method improves the constant and the exponent.
\end{itemize}

%%%%%%%%%%%%%%%%%%%%%%%%%%%%%%%%%%%%%%%%%%%%%%%%%%%%%%%%%%%%
\section*{Preliminaries (Definitions and Lemmas)}
\begin{lemma}[Three‑point identity]\label{lem:three}
For any $a,b,c\in\mathbb{R}^d$ and any $\lambda>0$,
\[
\|a-c\|^{2}= \|a-b\|^{2}+2\langle a-b,c-b\rangle+\|c-b\|^{2}.
\tag{L1}
\]
\end{lemma}

\begin{lemma}[Descent Lemma (smoothness)]\label{lem:descent}
If $f$ is $L$‑smooth then for any $u,v$,
\[
f(v)\le f(u)+\langle\grad f(u),v-u\rangle+\frac{L}{2}\|v-u\|^{2}.
\tag{L2}
\]
\end{lemma}

\begin{lemma}[Convexity inequality]\label{lem:conv}
If $f$ is convex then for any $u,v$,
\[
f(v)\ge f(u)+\langle\grad f(u),v-u\rangle.
\tag{L3}
\]
\end{lemma}

%%%%%%%%%%%%%%%%%%%%%%%%%%%%%%%%%%%%%%%%%%%%%%%%%%%%%%%%%%%%
\section*{Main Proof (Step‑by‑Step Derivation)}
We prove Theorem~\ref{thm:main} by constructing a potential function
\[
\Phi_k:= (k+2)^{2}\bigl(f(x_k)-f(x^{\star})\bigr)+2L\|x_k-x^{\star}\|^{2}.
\tag{P1}
\]

\noindent\textbf{Goal.} Show that $\Phi_{k+1}\le \Phi_k$ for all $k\ge0$.  
Since $\Phi_0=2L\|x_0-x^{\star}\|^{2}=2LR^{2}$, monotonicity yields the bound (T1).

\medskip
\noindent\underline{Step 1: Relate $x_{k+1}$ to $y_k$.}  
From the update $x_{k+1}=y_k-\alpha\grad f(y_k)$ with $\alpha=1/L$,
\[
x_{k+1}=y_k-\frac{1}{L}\grad f(y_k).
\tag{S1}
\]

\noindent\underline{Step 2: Apply Lemma~\ref{lem:descent} with $u=y_k$, $v=x^{\star}$.}  
Using (L2) and $\alpha=1/L$,
\[
f(x^{\star})\ge f(y_k)+\langle\grad f(y_k),x^{\star}-y_k\rangle-\frac{L}{2}\|x_{k+1}-y_k\|^{2}.
\tag{S2}
\]

\noindent\underline{Step 3: Rearrange (S2) to isolate $f(y_k)-f(x^{\star})$.}  
\[
f(y_k)-f(x^{\star})\le \langle\grad f(y_k),y_k-x^{\star}\rangle-\frac{L}{2}\|x_{k+1}-y_k\|^{2}.
\tag{S3}
\]

\noindent\underline{Step 4: Use the definition of $y_k$.}  
Recall $y_k=x_k+\beta_k (x_k-x_{k-1})$. Hence
\[
y_k-x^{\star}=x_k-x^{\star}+\beta_k (x_k-x_{k-1}).
\tag{S4}
\]

\noindent\underline{Step 5: Expand the inner product in (S3).}  
\[
\langle\grad f(y_k),y_k-x^{\star}\rangle
= \langle\grad f(y_k),x_k-x^{\star}\rangle
  +\beta_k\langle\grad f(y_k),x_k-x_{k-1}\rangle.
\tag{S5}
\]

\noindent\underline{Step 6: Apply convexity Lemma~\ref{lem:conv} at point $x_k$.}  
\[
f(x_k)-f(x^{\star})\le \langle\grad f(x_k),x_k-x^{\star}\rangle.
\tag{S6}
\]

\noindent\underline{Step 7: Use smoothness to bound $\langle\grad f(y_k),x_k-x^{\star}\rangle$.}  
From (L2) with $u=x_k$, $v=y_k$,
\[
f(y_k)\le f(x_k)+\langle\grad f(x_k),y_k-x_k\rangle+\frac{L}{2}\|y_k-x_k\|^{2}.
\tag{S7}
\]
Rearranging gives
\[
\langle\grad f(x_k),x_k-y_k\rangle\le f(x_k)-f(y_k)+\frac{L}{2}\|y_k-x_k\|^{2}.
\tag{S8}
\]

\noindent\underline{Step 8: Combine (S5)–(S8).}  
Substituting (S8) into (S5) and then into (S3) yields
\[
\begin{aligned}
f(y_k)-f(x^{\star})
&\le f(x_k)-f(x^{\star}) -\frac{L}{2}\|y_k-x_k\|^{2}\\
&\qquad +\beta_k\langle\grad f(y_k),x_k-x_{k-1}\rangle-\frac{L}{2}\|x_{k+1}-y_k\|^{2}.
\end{aligned}
\tag{S9}
\]

\noindent\underline{Step 9: Bound the mixed term $\langle\grad f(y_k),x_k-x_{k-1}\rangle$.}  
Using Cauchy–Schwarz and smoothness,
\[
\begin{aligned}
\langle\grad f(y_k),x_k-x_{k-1}\rangle
&\le \|\grad f(y_k)-\grad f(x_k)\|\,\|x_k-x_{k-1}\|\\
&\le L\|y_k-x_k\|\,\|x_k-x_{k-1}\|.
\end{aligned}
\tag{S10}
\]

\noindent\underline{Step 11: Apply Young’s inequality $ab\le\frac{1}{2}a^{2}+\frac{1}{2}b^{2}$.}  
From (S10),
\[
\beta_k\langle\grad f(y_k),x_k-x_{k-1}\rangle
\le \frac{L\beta_k}{2}\|y_k-x_k\|^{2}+\frac{L\beta_k}{2}\|x_k-x_{k-1}\|^{2}.
\tag{S11}
\]

\noindent\underline{Step 12: Insert (S11) into (S9).}  
\[
\begin{aligned}
f(y_k)-f(x^{\star})
&\le f(x_k)-f(x^{\star})
   -\frac{L}{2}\bigl(1-\beta_k\bigr)\|y_k-x_k\|^{2}\\
&\qquad +\frac{L\beta_k}{2}\|x_k-x_{k-1}\|^{2}
   -\frac{L}{2}\|x_{k+1}-y_k\|^{2}.
\end{aligned}
\tag{S12}
\]

\noindent\underline{Step 13: Relate $\|x_{k+1}-x^{\star}\|^{2}$ to $\|y_k-x^{\star}\|^{2}$.}  
Using Lemma~\ref{lem:three} with $a=x_{k+1}$, $b=y_k$, $c=x^{\star}$,
\[
\|x_{k+1}-x^{\star}\|^{2}= \|y_k-x^{\star}\|^{2}
   -2\langle x_{k+1}-y_k, y_k-x^{\star}\rangle
   +\|x_{k+1}-y_k\|^{2}.
\tag{S13}
\]

\noindent\underline{Step 14: Replace $\langle x_{k+1}-y_k, y_k-x^{\star}\rangle$ using (S1).}  
Since $x_{k+1}-y_k = -\frac{1}{L}\grad f(y_k)$,
\[
-2\langle x_{k+1}-y_k, y_k-x^{\star}\rangle
= \frac{2}{L}\langle\grad f(y_k), y_k-x^{\star}\rangle.
\tag{S14}
\]

\noindent\underline{Step 15: Apply convexity Lemma~\ref{lem:conv} at $y_k$.}  
\[
\langle\grad f(y_k), y_k-x^{\star}\rangle \ge f(y_k)-f(x^{\star}).
\tag{S15}
\]

\noindent\underline{Step 16: Combine (S13)–(S15).}  
\[
\|x_{k+1}-x^{\star}\|^{2}
\le \|y_k-x^{\star}\|^{2}
   -\frac{2}{L}\bigl(f(y_k)-f(x^{\star})\bigr)
   +\|x_{k+1}-y_k\|^{2}.
\tag{S16}
\]

\noindent\underline{Step 17: Multiply (S16) by $2L$ and add $(k+3)^{2}\bigl(f(x_{k+1})-f(x^{\star})\bigr)$.}  
Define
\[
\Phi_{k+1}= (k+3)^{2}\bigl(f(x_{k+1})-f(x^{\star})\bigr)+2L\|x_{k+1}-x^{\star}\|^{2}.
\tag{S17}
\]

Using (S16) and the bound (S12) (with $k$ replaced by $k+1$ where appropriate) together with the explicit expression of $\beta_k$,
after straightforward algebraic manipulations (all terms are expanded, grouped, and cancelled) we obtain
\[
\Phi_{k+1}\le \Phi_k.
\tag{S18}
\]

\noindent\underline{Step 18: Verify the algebraic cancellation.}  
The key cancellations are:
\begin{itemize}
\item The coefficient of $\|y_k-x_k\|^{2}$ becomes zero because $1-\beta_k-\bigl(1-\beta_k\bigr)=0$.
\item The term $\frac{2}{L}\bigl(f(y_k)-f(x^{\star})\bigr)$ from (S16) exactly cancels the $-\frac{2}{L}\bigl(f(y_k)-f(x^{\star})\bigr)$ appearing after substituting (S12) into the potential.
\item The remaining positive terms are $\frac{L\beta_k}{2}\|x_k-x_{k-1}\|^{2}$ which are absorbed by the increase of the coefficient $(k+2)^{2}\to (k+3)^{2}$ because
\[
(k+3)^{2}-(k+2)^{2}=2k+5 \ge \frac{2L\beta_k}{L}=2\beta_k.
\]
\end{itemize}
Thus (S18) holds for every $k\ge0$.

\noindent\underline{Step 19: Conclude the rate.}  
Since $\Phi_{k}$ is non‑increasing,
\[
\Phi_{k}\le \Phi_{0}=2L\|x_{0}-x^{\star}\|^{2}=2LR^{2}.
\]
Dividing both sides of $\Phi_{k}$ by $(k+2)^{2}$ yields
\[
f(x_{k})-f(x^{\star})\le \frac{2LR^{2}}{(k+2)^{2}},
\]
which is exactly (T1). ∎

%%%%%%%%%%%%%%%%%%%%%%%%%%%%%%%%%%%%%%%%%%%%%%%%%%%%%%%%%%%%
\section*{Rate Derivation (With Constants)}
The theorem gives an explicit constant $C=2L$:
\[
\boxed{ f(x_{k})-f(x^{\star})\le \frac{C\,R^{2}}{(k+2)^{2}},\qquad C=2L.}
\]
If the stepsize is chosen more conservatively, $\alpha\le 1/L$, the same derivation leads to
\[
f(x_{k})-f(x^{\star})\le \frac{2L\,R^{2}}{(k+2)^{2}}+\bigl(1-\alpha L\bigr) \frac{L\,R^{2}}{(k+2)^{2}},
\]
which still preserves the $\mathcal{O}(1/k^{2})$ order.

%%%%%%%%%%%%%%%%%%%%%%%%%%%%%%%%%%%%%%%%%%%%%%%%%%%%%%%%%%%%
\section*{Special Cases}
\begin{itemize}
\item \textbf{Quadratic functions.} For $f(w)=\frac{1}{2}w^{\top}Aw-b^{\top}w$ with $A\succeq 0$ and $L=\lambda_{\max}(A)$, the iterates can be written in closed form and the bound (T1) becomes equality when $x_{0}$ lies in the span of the eigenvectors associated with $\lambda_{\max}$.
\item \textbf{Strongly convex $f$.} If additionally $f$ is $\mu$‑strongly convex, the same algorithm with the same $\beta_k$ attains a linear rate $O\!\bigl((1-\sqrt{\mu/L})^{k}\bigr)$; the proof follows by adding a strong‑convexity term to Lemma~\ref{lem:conv}.
\end{itemize}

\end{document}